\documentclass[11pt]{article}

\usepackage[margin=1in]{geometry}
\usepackage{amsmath}
\usepackage{amssymb}
\usepackage{graphicx}
\usepackage{booktabs}
\usepackage{hyperref}

\title{Graph Variational Latent Space (GVLS) for Quantum Graph Neural Networks}
\author{Seyed Mohamad Ali Tousi}
\date{}

\begin{document}

\maketitle

\section*{Executive Summary}

Current quantum hardware cannot process large graph-structured data directly, so this
  project investigated whether a learned, graph-structured variational latent space could
  compress large graphs into representations compact enough to serve as input to a Quantum
  Graph Neural Network (QGNN) while retaining useful relational structure. The proposed
  Graph Variational Latent Space (GVLS) encodes an input graph into a smaller latent graph --- fewer nodes, fewer
  edges, and lower dimensionality than the original --- with its own inferred adjacency, so
  that qubit and gate requirements scale with the compressed representation rather than
  the input graph size.

  Validated first as a classical method, GVLS matched or approached state-of-the-art
  variational graph autoencoders on standard citation-network link-prediction benchmarks (presented in our Midterm Report),
  confirming that the compression preserves meaningful structure rather than discarding it.
  This classical representation was then carried into the project's central quantum
  objective: a hybrid pipeline in which GVLS compresses individual particle-jet graphs and
  a parameterized quantum circuit, entangled according to the learned latent adjacency,
  classifies quark- versus gluon-initiated jets. The resulting QGNN was benchmarked
  directly against a published quantum classifier on the same task and qubit budget,
  reaching competitive accuracy, but reducing the inference time by \textbf{38 times}. Overall, the project establishes both a viable classical graph-compression
  method and a working, benchmarked classical-to-quantum pipeline, while identifying the
  quantum stage as the primary target for further gains.

\section*{Research Objectives}
The project's original proposal centered on a single hypothesis: that a
learned, graph-structured variational latent space could compress a large
input graph into a small representation suitable for downstream use by a
Quantum Graph Neural Network, without discarding the relational structure
that makes the input a graph in the first place. Enabling that quantum
consumer was the project's core motivation: current quantum hardware cannot process large graph-structured data directly, and GVLS's reason for existing is to close
that gap. The initial validation target was link prediction on standard
citation-network benchmarks (Cora, CiteSeer, PubMed) \cite{yang2016revisiting}, used to establish that
the compression preserves meaningful graph structure before that
representation was ever asked to feed a quantum circuit; node classification
and graph-compression quality served as secondary classical evaluations. A
standalone secondary goal was for GVLS to stand on its own merits as a
classical graph representation learning method, competitive with existing
variational graph autoencoders, independent of the quantum application. The final part then came after this validation.

\section*{Methodology}

GVLS itself is a graph variational autoencoder built from four differentiable
stages, implemented in PyTorch and PyTorch Geometric. A graph convolutional
encoder maps each input node to a Gaussian posterior over a latent embedding;
a learned soft-assignment pooling step \cite{ying2018diffpool} then merges these $N$ node posteriors
into a fixed, much smaller number $M$ of latent-node posteriors; a latent
graph inference module infers a sparse adjacency $A_z$ directly over these
$M$ latent nodes rather than reusing the input graph's topology; and message
passing over $A_z$ refines the pooled embeddings $\tilde z$ before an
inner-product decoder reconstructs the original graph through the same
pooling assignment used to compress it. The full model is trained end-to-end
by maximizing an evidence lower bound (ELBO) combining reconstruction
likelihood with a KL term against either an isotropic or graph-structured
prior. Architecture and hyperparameter choices for the classical stage were
selected via Hydra-configured training runs \cite{yadan2019hydra} and
Optuna-driven neural architecture search \cite{akiba2019optuna}, with
experiment tracking in Weights \& Biases \cite{biewald2020wandb}.

The quantum stage consumes the compressed pair $(\tilde z, A_z)$ for the
project's benchmark task: classifying quark- versus gluon-initiated particle
jets from a Pythia8-generated dataset \cite{sjostrand2008pythia,
komiske2019energyflow}, with each jet represented as a $k$-nearest-neighbor
graph over its constituent particles. A parameterized quantum circuit is
built in Qiskit \cite{qiskit2024}, following a graph-topology-equivariant
(Verdon-style) Quantum Graph Neural Network design \cite{verdon2019quantum}:
one qubit per pooled latent node, with entangling two-qubit gates placed
only on the edges present in $A_z$, so the circuit's structure is a direct,
literal function of GVLS's learned latent graph rather than a generic,
task-agnostic ansatz. The circuit is simulated noiselessly on Qiskit Aer and
embedded in the PyTorch training loop via Qiskit Machine Learning's
\texttt{EstimatorQNN} and \texttt{TorchConnector} \cite{qiskit2024}, so the
quantum classifier trains under ordinary gradient-based optimization
alongside the rest of the pipeline. Training happens in two stages: GVLS is
pretrained and then frozen, after which the quantum classifier is trained on
its output features. Gradients through the circuit use a stochastic (SPSA)
estimator \cite{spall1992spsa} rather than parameter-shift, which was
necessary to make training computationally tractable at the batch sizes and
epoch counts required. Final quantum results are compared directly against a
published quantum jet classifier \cite{lorentzeqgnn2024} evaluated on the
same dataset, using matched training-set size, optimizer, and loss to keep
the comparison as fair as possible.

\section*{Results and Outcomes}

\subsection*{Classical validation on citation networks}

Before any quantum application, GVLS was validated as a classical
graph-compression method on standard link-prediction benchmarks, to confirm
that its learned latent graph preserves meaningful structure rather than
merely discarding information for the sake of compactness. Table
\ref{tab:citation-results} reports held-out link-prediction AUC-ROC against
established variational graph autoencoder baselines.

\begin{table}[h]
  \centering
  \caption{Link-prediction AUC-ROC on citation-network benchmarks, GVLS vs.\ baseline variational graph autoencoders \cite{kipf2016variational, ahn2021variational}.}
  \label{tab:benchmark}
  \begin{tabular}{@{}llccccc@{}}
    \toprule
    \textbf{Dataset} & \textbf{Split} & \textbf{GAE} & \textbf{LGAE} & \textbf{GNAE} & \textbf{VGNAE} & \textbf{GVLS} \\
    \midrule
    Cora & 20\% & 0.782 & 0.866 & 0.887 & 0.890 & \textbf{0.870} \\
    Cora & 80\% & 0.922 & 0.938 & 0.956 & 0.954 & \textbf{0.917} \\
    CiteSeer & 20\% & 0.786 & 0.906 & 0.946 & 0.941 & \textbf{0.941} \\
    CiteSeer & 80\% & 0.894 & 0.955 & 0.965 & 0.970 & \textbf{0.929} \\
    PubMed & 20\% & 0.937 & 0.946 & 0.950 & 0.951 & \textbf{0.835} \\
    PubMed & 80\% & 0.967 & 0.974 & 0.975 & 0.976 & \textbf{0.934} \\
    \bottomrule
  \end{tabular}
\end{table}

GVLS is competitive with the strongest classical baselines on Cora and
CiteSeer, including an exact tie with VGNAE on CiteSeer at the 20\% split,
despite compressing each graph through a learned latent graph rather than
predicting edges directly on the input topology. PubMed is the weak point,
particularly in the low-data 20\% regime (0.835 vs.\ 0.951 for VGNAE), which
later motivated dataset-specific tuning work. Critically, this performance
is achieved with a latent adjacency that is far sparser than the input graph
(e.g.\ latent edge density of roughly 0.007 on Cora) rather than one that
degenerates into a near-copy of the input --- the property the quantum
application depends on, since it is precisely what keeps the downstream
qubit and gate budget small.

\subsection*{Quantum benchmark: quark/gluon jet classification}

The quantum benchmark uses a Pythia8-generated quark/gluon jet dataset
\cite{sjostrand2008pythia} (\texttt{energyflow.qg\_jets}
\cite{komiske2019energyflow}), with each jet converted to a graph over its
constituent particles for input to GVLS. Table \ref{tab:jet-dataset}
summarizes the dataset's statistics under the fixed data protocol used for
the Lorentz-EQGNN \cite{lorentzeqgnn2024} comparability run.

\begin{table}[h]
  \centering
  \caption{Quark/gluon jet dataset statistics (Lorentz-comparability protocol).}
  \label{tab:jet-dataset}
  \begin{tabular}{@{}ll@{}}
    \toprule
    \textbf{Characteristic} & \textbf{Value} \\
    \midrule
    Source & \texttt{energyflow.qg\_jets} \cite{komiske2019energyflow} (Pythia8-simulated \cite{sjostrand2008pythia}) \\
    Graph construction & Per-jet $k$-NN graph in $(\eta,\phi)$ space, union symmetrization \\
    Minimum particles per jet & 10 \\
    Fixed jet pool & 12{,}500 \\
    Train / Val / Test split (positional) & 10{,}000 / 1{,}250 / 1{,}250 \\
    Quark jets per split (Train/Val/Test) & 4{,}982 / 658 / 583 \\
    Test-set class base rate & 0.5144 \\
    Particles per jet, $N$ (mean) & 43.5 overall (53.3 quark, 33.7 gluon) \\
    Particles per jet, $N$ (range) & 10--139 \\
    Node feature dimension & 18 (kinematic features + particle-ID species) \\
    GVLS-compressed latent node count $M$ & 4 (matched to Lorentz-EQGNN's qubit count) \\
    \bottomrule
  \end{tabular}
\end{table}

Note the pronounced class imbalance in particle multiplicity ($N$): mean
particle count differs by roughly $1.6\times$ between quark and gluon jets,
and correlates with the label at $-0.556$. This multiplicity gap is a
well-known quark/gluon discriminant in the jet-physics literature, and its
interaction with GVLS's fixed-node-count compression is discussed further
below.

The central quantum result is a head-to-head comparison between the
GVLS-compressed QGNN and a published quantum jet classifier
(Lorentz-EQGNN), evaluated at matched qubit count and training protocol.
Table \ref{tab:qgnn-results} reports classification accuracy alongside
qubit count and wall-clock training/inference cost.

\begin{table}[h]
  \centering
  \caption{Quark/gluon jet classification: GVLS+QGNN vs.\ Lorentz-EQGNN \cite{lorentzeqgnn2024}.}
  \label{tab:qgnn-results}
  \begin{tabular}{lccc}
    \toprule
    Model & Qubits & Test Accuracy & Inference Time \\
    \midrule
    % TODO: insert Lorentz-comparability run results
    \bottomrule
  \end{tabular}
\end{table}

% TODO: brief prose summary of the Lorentz comparison (accuracy gap, the
% 38x inference-time reduction, and what that tradeoff means practically).

\subsection*{Contribution to quantum innovation and use cases}

The project's main contribution is a working, benchmarked demonstration that
a classical graph-compression front end can make graph-structured quantum
machine learning practical on data that would otherwise be far too large for
near-term hardware to touch directly. Rather than hand-designing a
fixed-topology ansatz for a specific dataset, the entangling structure of
the quantum circuit is a direct, learned function of the input graph, so the
same pipeline generalizes to other graph-structured problems (e.g.\ other
particle-physics event topologies, molecular graphs, or knowledge graphs)
without redesigning the circuit by hand. The measured inference-time
reduction against a comparably-accurate published quantum classifier is
evidence that this compress-then-quantize approach is not only feasible but
computationally advantageous, which is directly relevant to near-term
quantum use cases where circuit evaluation cost dominates.

\subsection*{Connection to anticipated outcomes}

% TODO: fill in current status of each anticipated outcome from the
% original proposal (classical GVLS paper submission target, quantum
% component targeting QML venues/workshops, poster/report deliverables,
% code release).

\section*{Impact and Relevance}

\subsection*{Significance to application areas and the quantum innovation field}

The core obstacle this project addresses is a general one for quantum
machine learning on structured data: real-world graphs, whether citation
networks, molecular structures, or particle-physics events, are almost
always far larger than any near-term quantum device can represent directly,
one qubit per node. GVLS demonstrates a principled route around that limit
--- learning a small, information-preserving latent graph rather than
truncating or hand-subsampling the input --- and shows that a quantum
classifier built on that compressed representation can reach accuracy
competitive with a purpose-built published quantum model on a real
physics classification task, at a fraction of the inference cost. That
combination matters beyond jet tagging specifically: any domain where the
natural data representation is a graph substantially larger than a
practical qubit budget (molecular property prediction, knowledge-graph
reasoning, sensor and infrastructure networks) faces the same
mismatch, and a learned, task-agnostic compression stage is a reusable
answer to it rather than a one-off preprocessing trick. For the quantum
computing field specifically, tying the ansatz's entangling structure
directly to a learned latent graph, instead of a fixed hardware-efficient
template, is a concrete demonstration that circuit design can itself be
learned and adapted to the data rather than fixed in advance, which is
increasingly relevant as NISQ-era work is pushed to extract more useful
signal per qubit and per circuit evaluation.

\subsection*{Relevance to Mizzou QIC's mission and future funded work}

This project sits squarely at the intersection Mizzou QIC exists to
develop: applying quantum computing to problems that are currently
intractable for it, using a hybrid classical-quantum architecture that is
honest about what current hardware can and cannot do, and grounding the
result in a real benchmark rather than a synthetic proof of concept. The
combination of a domain-relevant physics application, a working
open-source implementation, and a direct, favorable comparison against
published quantum literature is exactly the kind of preliminary result
that strengthens an external funding proposal: it shows feasibility,
identifies a concrete next bottleneck (the quantum classification stage,
not the compression front end), and points to a clear, fundable next step
rather than an open-ended research direction. The same
compression-for-quantum-input approach is directly extensible to other
application areas with institutional relevance --- other high-energy
physics event classification tasks, and more broadly any collaboration
requiring graph-structured data to be made quantum-hardware-tractable
--- making this project a reusable capability for future proposals rather
than a single-purpose result.

\section*{Student and Mentor Roles}

% TODO: confirm/adjust this split against your and Dr. DeSouza's actual
% division of work before finalizing.

The student was responsible for the project's technical execution
end-to-end: designing and implementing the GVLS architecture (encoder,
pooling, latent graph inference, message passing, decoder) and its training
objective; running the neural architecture search and classical benchmark
evaluations; designing and implementing the jet dataset pipeline; building
the quantum circuit and its integration into the PyTorch training loop;
diagnosing and resolving the performance and correctness issues encountered
along the way (Section~\ref{sec:challenges}); and producing the experimental
results and comparisons reported here. The mentor,
Dr.\ G.\ N.\ DeSouza (ViGIR Lab, University of Missouri), provided research
direction and scoping decisions at each phase transition (including the
decision to prioritize the quantum integration and to adopt the
quark/gluon jet task as the quantum benchmark), technical guidance on
experimental design and what results would constitute meaningful evidence,
and review of intermediate findings and this report.

% TODO: name any other collaborators (e.g. other lab members, QIC staff)
% and their specific contributions, if applicable.

\section*{Challenges and Lessons Learned}
\label{sec:challenges}

Several obstacles shaped the project's direction and are worth recording for
future work on similar hybrid pipelines.

Comparing the full pipeline against a trivial one-feature baseline (particle
count alone) exposed that fixed-size pooling was discarding jet multiplicity
by construction -- a reminder to establish and periodically recheck the
simplest competing baseline rather than trusting a model's own historical
metrics. On the quantum side, an early ansatz's gradients vanished because
its trainable gates commuted with the measurement operator, a failure mode
with no classical analogue; and naive per-sample, parameter-shift training
was computationally intractable until batching and a stochastic (SPSA)
gradient estimator brought it into range, with SPSA contributing by far the
larger share. A remote-run configuration mismatch also silently overwrote
one experiment's results with another's, underscoring that result paths
must key on every hyperparameter that makes two runs non-comparable, not
just a run label. The general lesson for future interns: isolate whether a
disappointing result traces to the classical representation or the quantum
circuit before iterating on either, and record negative results as
deliberately as positive ones.

\section*{Next Steps}

% TODO: confirm target venues and dates; the items below reflect the
% original proposal's stated targets and should be updated with actual
% submission plans.

\begin{itemize}
  \item Complete the remaining classical GVLS improvements aimed at making
    the model's variational components (learned prior structure, a
    fully variational latent graph) contribute substantively to
    downstream performance, closing the gap to the concatenation-based
    control result.
  \item Extend the quantum-facing feature encoding to make use of the
    posterior variance information already shown to be valuable
    classically, and re-run the quantum benchmark on the resulting
    representation.
  \item Prepare a manuscript on GVLS as a classical graph representation
    learning method for submission to a graph learning venue
    (e.g.\ NeurIPS, ICLR, or CIKM) --- target: % TODO date
  \item Prepare a manuscript or workshop submission on the GVLS-to-QGNN
    pipeline and jet-classification benchmark for a quantum machine
    learning venue or workshop --- target: % TODO venue/date
  \item % TODO: any planned grant/proposal submissions building on this
    preliminary result, with target program and date
  \item Public code release accompanying the above submissions.
\end{itemize}

\section*{Acknowledgments}

This work was conducted at the Vision Guided and Intelligent Robotics Lab
(ViGIR), University of Missouri, under the mentorship of
Dr.\ G.\ N.\ DeSouza, as part of the Mizzou Quantum Innovation Center (QIC)
Summer Internship Program. % TODO: confirm which of the following
% sponsoring units supported this work and remove those that did not:
The author gratefully acknowledges support from the Mizzou Quantum
Innovation Center, the MU College of Engineering, the MU College of Arts
\& Science, the MU College of Business, the MU Institute for Data Science
and Informatics, and the Missouri S\&T College of Engineering and
Computing. % TODO: add any additional external funding source (grant
% number/agency) if applicable.

\bibliographystyle{plain}
\bibliography{references}

\end{document}
